Tabulka \ref{tab:coverage_correlation_table} shrnuje výsledky korelační analýzy provedené na párovém panelu zemí OECD/\ac{EU} za dostupná roky od roku 2004.
Zdrojem dat pro pokrytí kolektivního vyjednávání je databáze OECD/AIAS ICTWSS (ukazatel CBC, míra ERB), pro hustotu odborů tatáž databáze (TUD); výstupní proměnné pocházejí z~Eurostatu (lfsa\_ewhun2, nama\_10\_pc, ilc\_di12, earn\_nt\_net).
Každý řádek reportuje: počet párových pozorování $n$ (geo$\times$rok po párování datových sad), Pearsonovo $r$ na celém panelu (míra lineární asociace), Spearmanovo $\rho$ jako robustní neparametrická alternativa rezistentní vůči odlehlým hodnotám, a~$r_{\text{akt.}}$ = Pearsonovo $r$ pro průřezová data z~roku 2024.
Síla korelací je hodnocena dle Evansovy (1996) škály: $|r| < 0{,}20$ velmi slabá; $0{,}20$--$0{,}39$ slabá; $0{,}40$--$0{,}59$ střední; $0{,}60$--$0{,}79$ silná; $\geq 0{,}80$ velmi silná.

Výsledky jsou interpretovány po jednotlivých párech:

\begin{itemize}
  \item \textbf{Pokrytí \ac{KV} vs.~pracovní doba} ($r = -0{,}492$; $\rho = -0{,}384$):
    panelová Pearsonova korelace svědčí o~\emph{střední záporné} asociaci --- vyšší pokrytí \ac{KS} je systematicky spojeno s~kratší průměrnou pracovní dobou.
    Mírný pokles na Spearmanovo $\rho$ (slabá--střední) naznačuje přítomnost několika málo vlivných pozorování posouvajících Pearsonův koeficient, nikoliv výrazný outlier narušující charakter vztahu.
    Průřezová korelace $r_{\text{akt.}} = -0{,}536$ (\emph{střední záporná}) se oproti panelové hodnotě spíše zesílila --- vztah tedy není artefaktem historické variace, ale přetrvává i~v~aktuálním mezinárodním průřezu.

  \item \textbf{Pokrytí \ac{KV} vs.~\ac{HDP}/ob.~(\ac{PPS})} ($r = 0{,}471$; $\rho = 0{,}456$):
    \emph{střední kladná} asociace, vysoce konzistentní mezi Pearsonem a~Spearmanem --- rozdíl pouhých 0,015 indikuje přibližně lineární vztah bez dominantních odlehlých pozorování.
    Průřezová hodnota $r_{\text{akt.}} = 0{,}598$ se přibližuje hranici \emph{silné korelace} (Evans: 0,60), naznačujíce, že v~aktuálním průřezu roku 2024 je vazba mezi institucemi \ac{KV} a~životní úrovní ještě zřetelnější než v~historickém panelu.

  \item \textbf{Pokrytí \ac{KV} vs.~Gini} ($r = -0{,}445$; $\rho = -0{,}286$):
    tento pár vykazuje nejmarkantnější divergenci mezi Pearsonem a~Spearmanem (rozdíl 0,159), která diagnosticky ukazuje na přítomnost outlierů zvyšujících Pearsonův koeficient.
    Nejpravděpodobnějším kandidátem je Francie: s~pokrytím \ac{KS} přes \SI{95}{\percent} (erga omnes) vykazuje relativně vysoké Gini (30--32), protože smluvní komprese mezd je přebita vysokou nerovností kapitálových příjmů a~vrcholných odměn.
    Průřezová hodnota $r_{\text{akt.}} = -0{,}238$ (\emph{slabá záporná}) dále ilustruje, proč nelze z~jednoduchého průřezu soudit o~teorii: v~roce 2024 komprimují Gini napříč EU27 silné transferové systémy a~progresivní zdanění, přičemž \ac{KV} je jen jedním z~více kolineárních institucio\-nálních faktorů.
    Panelový koeficient ($-0{,}445$) sledující \emph{změny} v~čase v~rámci zemí je proto informačně hodnotnější: dokumentuje, že kdy\-koli pokrytí \ac{KS} v~dané zemi kleslo, Gini se tendovalo zvýšit.

  \item \textbf{Pokrytí \ac{KV} vs.~příjem/\ac{HDP}} ($r = 0{,}227$; $\rho = 0{,}230$):
    \emph{slabá kladná} asociace, numericky konzistentní v~obou metodách.
    Relativně nízká hodnota je z~velké části artefaktem konstrukce ukazatele: \ac{HDP}/ob.~je sám o~sobě pozitivně korelován s~\ac{KV} pokrytím ($r = 0{,}471$), takže ve jmenovateli podílu čistý~příjem/\ac{HDP} oba členy rostou společně s~pokrytím a~efekt se vzájemně tlumí.
    Tento výsledek tedy nesvědčí o~tom, že \ac{KV} instituce nezvedají mzdy, ale o~tom, že tento konkrétní relativní ukazatel není citlivým testem dané hypotézy --- absolutní příjmová míra (předchozí pár) zůstává spolehlivějším indikátorem.

  \item \textbf{Hustota odborů vs.~Gini} ($r = -0{,}454$; $\rho = -0{,}455$):
    tento pár je statisticky nejrobustnějším v~celé tabulce: Pearsonovo $r$ a~Spearmanovo $\rho$ jsou prakticky totožné (rozdíl $< 0{,}001$), což svědčí o~přibližně lineárním vztahu bez systematického vlivu odlehlých hodnot.
    \emph{Střední záporná} asociace je stabilní v~celém panelu i~v~aktuálním průřezu ($r_{\text{akt.}} = -0{,}541$).
    Shoda obou koeficientů rovněž zpětně potvrzuje, že divergence u~páru \ac{KV}~vs.~Gini je specifická pro pokrytí (kde Francie tvoří strukturální outlier), nikoliv obecná vlastnost vztahu institucí \ac{KV} k~nerovnosti.
\end{itemize}

Žádná z~výše uvedených korelací nefalzifikuje teorii mzdové komprese prostřednictvím kolektivního vyjednávání.
Klíčovým argumentem je, že jednoduchou bivariátní korelací na heterogenním průřezu 27 zemí lze jen obtížně izolovat marginální příspěvek \ac{KV} pokrytí --- vzdor tomu jsou všechny korelace ve správném předpokládaném směru bez jediné výjimky.
Kauzální identifikaci přináší panelová analýza s~fixními efekty a~přirozené experimenty: metaanalytická literatura (např.~Jaumotte \& Osorio Buitron 2015; Card et al.~2004) konzistentně dokládá statisticky i~věcně významný efekt \ac{KV} institucí na mzdovou kompresi.
Korelace v~tabulce tedy plní funkci deskriptivního podkladu motivujícího institucionální případ\-ovou studii ČR --- nikoliv jako izolovaný kauzální důkaz, ale jako konzistentní mezinárodní vzorec, v~němž ČR se svým pokrytím $\approx 32\,\%$ systematicky vykazuje horšíu výsledky než srovnatelné ekonomiky s~vyšším pokrytím \ac{KS}.

Tabulka \ref{tab:coverage_correlation_table} shrnuje výsledky korelační analýzy provedené na párovém panelu zemí OECD/\ac{EU} za dostupná roky od roku 2004.
Zdrojem dat pro pokrytí kolektivního vyjednávání je databáze OECD/AIAS ICTWSS (ukazatel CBC, míra ERB), pro hustotu odborů tatáž databáze (TUD); výstupní proměnné pocházejí z~Eurostatu (lfsa\_ewhun2, nama\_10\_pc, ilc\_di12, earn\_nt\_net).
Každý řádek reportuje: počet párových pozorování $n$ (geo$\times$rok po párování datových sad), Pearsonovo $r$ na celém panelu (míra lineární asociace), Spearmanovo $\rho$ jako robustní neparametrická alternativa rezistentní vůči odlehlým hodnotám, a~$r_{\text{akt.}}$ = Pearsonovo $r$ pro průřezová data z~roku 2024.
Síla korelací je hodnocena dle Evansovy (1996) škály: $|r| < 0{,}20$ velmi slabá; $0{,}20$--$0{,}39$ slabá; $0{,}40$--$0{,}59$ střední; $0{,}60$--$0{,}79$ silná; $\geq 0{,}80$ velmi silná.

Výsledky jsou interpretovány po jednotlivých párech:

\begin{itemize}
  \item \textbf{Pokrytí \ac{KV} vs.~pracovní doba} ($r = -0{,}492$; $\rho = -0{,}384$):
    panelová Pearsonova korelace svědčí o~\emph{střední záporné} asociaci --- vyšší pokrytí \ac{KS} je systematicky spojeno s~kratší průměrnou pracovní dobou.
    Mírný pokles na Spearmanovo $\rho$ (slabá--střední) naznačuje přítomnost několika málo vlivných pozorování posouvajících Pearsonův koeficient, nikoliv výrazný outlier narušující charakter vztahu.
    Průřezová korelace $r_{\text{akt.}} = -0{,}536$ (\emph{střední záporná}) se oproti panelové hodnotě spíše zesílila --- vztah tedy není artefaktem historické variace, ale přetrvává i~v~aktuálním mezinárodním průřezu.

  \item \textbf{Pokrytí \ac{KV} vs.~\ac{HDP}/ob.~(\ac{PPS})} ($r = 0{,}471$; $\rho = 0{,}456$):
    \emph{střední kladná} asociace, vysoce konzistentní mezi Pearsonem a~Spearmanem --- rozdíl pouhých 0,015 indikuje přibližně lineární vztah bez dominantních odlehlých pozorování.
    Průřezová hodnota $r_{\text{akt.}} = 0{,}598$ se přibližuje hranici \emph{silné korelace} (Evans: 0,60), naznačujíce, že v~aktuálním průřezu roku 2024 je vazba mezi institucemi \ac{KV} a~životní úrovní ještě zřetelnější než v~historickém panelu.

  \item \textbf{Pokrytí \ac{KV} vs.~Gini} ($r = -0{,}445$; $\rho = -0{,}286$):
    tento pár vykazuje nejmarkantnější divergenci mezi Pearsonem a~Spearmanem (rozdíl 0,159), která diagnosticky ukazuje na přítomnost outlierů zvyšujících Pearsonův koeficient.
    Nejpravděpodobnějším kandidátem je Francie: s~pokrytím \ac{KS} přes \SI{95}{\percent} (erga omnes) vykazuje relativně vysoké Gini (30--32), protože smluvní komprese mezd je přebita vysokou nerovností kapitálových příjmů a~vrcholných odměn.
    Průřezová hodnota $r_{\text{akt.}} = -0{,}238$ (\emph{slabá záporná}) dále ilustruje, proč nelze z~jednoduchého průřezu soudit o~teorii: v~roce 2024 komprimují Gini napříč EU27 silné transferové systémy a~progresivní zdanění, přičemž \ac{KV} je jen jedním z~více kolineárních institucio\-nálních faktorů.
    Panelový koeficient ($-0{,}445$) sledující \emph{změny} v~čase v~rámci zemí je proto informačně hodnotnější: dokumentuje, že kdy\-koli pokrytí \ac{KS} v~dané zemi kleslo, Gini se tendovalo zvýšit.

  \item \textbf{Pokrytí \ac{KV} vs.~příjem/\ac{HDP}} ($r = 0{,}227$; $\rho = 0{,}230$):
    \emph{slabá kladná} asociace, numericky konzistentní v~obou metodách.
    Relativně nízká hodnota je z~velké části artefaktem konstrukce ukazatele: \ac{HDP}/ob.~je sám o~sobě pozitivně korelován s~\ac{KV} pokrytím ($r = 0{,}471$), takže ve jmenovateli podílu čistý~příjem/\ac{HDP} oba členy rostou společně s~pokrytím a~efekt se vzájemně tlumí.
    Tento výsledek tedy nesvědčí o~tom, že \ac{KV} instituce nezvedají mzdy, ale o~tom, že tento konkrétní relativní ukazatel není citlivým testem dané hypotézy --- absolutní příjmová míra (předchozí pár) zůstává spolehlivějším indikátorem.

  \item \textbf{Hustota odborů vs.~Gini} ($r = -0{,}454$; $\rho = -0{,}455$):
    tento pár je statisticky nejrobustnějším v~celé tabulce: Pearsonovo $r$ a~Spearmanovo $\rho$ jsou prakticky totožné (rozdíl $< 0{,}001$), což svědčí o~přibližně lineárním vztahu bez systematického vlivu odlehlých hodnot.
    \emph{Střední záporná} asociace je stabilní v~celém panelu i~v~aktuálním průřezu ($r_{\text{akt.}} = -0{,}541$).
    Shoda obou koeficientů rovněž zpětně potvrzuje, že divergence u~páru \ac{KV}~vs.~Gini je specifická pro pokrytí (kde Francie tvoří strukturální outlier), nikoliv obecná vlastnost vztahu institucí \ac{KV} k~nerovnosti.
\end{itemize}

Žádná z~výše uvedených korelací nefalzifikuje teorii mzdové komprese prostřednictvím kolektivního vyjednávání.
Klíčovým argumentem je, že jednoduchou bivariátní korelací na heterogenním průřezu 27 zemí lze jen obtížně izolovat marginální příspěvek \ac{KV} pokrytí --- vzdor tomu jsou všechny korelace ve správném předpokládaném směru bez jediné výjimky.
Kauzální identifikaci přináší panelová analýza s~fixními efekty a~přirozené experimenty: metaanalytická literatura (např.~Jaumotte \& Osorio Buitron 2015; Card et al.~2004) konzistentně dokládá statisticky i~věcně významný efekt \ac{KV} institucí na mzdovou kompresi.
Korelace v~tabulce tedy plní funkci deskriptivního podkladu motivujícího institucionální případ\-ovou studii ČR --- nikoliv jako izolovaný kauzální důkaz, ale jako konzistentní mezinárodní vzorec, v~němž ČR se svým pokrytím $\approx 32\,\%$ systematicky vykazuje horšíu výsledky než srovnatelné ekonomiky s~vyšším pokrytím \ac{KS}.

Tabulka \ref{tab:coverage_correlation_table} shrnuje výsledky korelační analýzy provedené na párovém panelu zemí OECD/\ac{EU} za dostupná roky od roku 2004.
Zdrojem dat pro pokrytí kolektivního vyjednávání je databáze OECD/AIAS ICTWSS (ukazatel CBC, míra ERB), pro hustotu odborů tatáž databáze (TUD); výstupní proměnné pocházejí z~Eurostatu (lfsa\_ewhun2, nama\_10\_pc, ilc\_di12, earn\_nt\_net).
Každý řádek reportuje: počet párových pozorování $n$ (geo$\times$rok po párování datových sad), Pearsonovo $r$ na celém panelu (míra lineární asociace), Spearmanovo $\rho$ jako robustní neparametrická alternativa rezistentní vůči odlehlým hodnotám, a~$r_{\text{akt.}}$ = Pearsonovo $r$ pro průřezová data z~roku 2024.
Síla korelací je hodnocena dle Evansovy (1996) škály: $|r| < 0{,}20$ velmi slabá; $0{,}20$--$0{,}39$ slabá; $0{,}40$--$0{,}59$ střední; $0{,}60$--$0{,}79$ silná; $\geq 0{,}80$ velmi silná.

Výsledky jsou interpretovány po jednotlivých párech:

\begin{itemize}
  \item \textbf{Pokrytí \ac{KV} vs.~pracovní doba} ($r = -0{,}492$; $\rho = -0{,}384$):
    panelová Pearsonova korelace svědčí o~\emph{střední záporné} asociaci --- vyšší pokrytí \ac{KS} je systematicky spojeno s~kratší průměrnou pracovní dobou.
    Mírný pokles na Spearmanovo $\rho$ (slabá--střední) naznačuje přítomnost několika málo vlivných pozorování posouvajících Pearsonův koeficient, nikoliv výrazný outlier narušující charakter vztahu.
    Průřezová korelace $r_{\text{akt.}} = -0{,}536$ (\emph{střední záporná}) se oproti panelové hodnotě spíše zesílila --- vztah tedy není artefaktem historické variace, ale přetrvává i~v~aktuálním mezinárodním průřezu.

  \item \textbf{Pokrytí \ac{KV} vs.~\ac{HDP}/ob.~(\ac{PPS})} ($r = 0{,}471$; $\rho = 0{,}456$):
    \emph{střední kladná} asociace, vysoce konzistentní mezi Pearsonem a~Spearmanem --- rozdíl pouhých 0,015 indikuje přibližně lineární vztah bez dominantních odlehlých pozorování.
    Průřezová hodnota $r_{\text{akt.}} = 0{,}598$ se přibližuje hranici \emph{silné korelace} (Evans: 0,60), naznačujíce, že v~aktuálním průřezu roku 2024 je vazba mezi institucemi \ac{KV} a~životní úrovní ještě zřetelnější než v~historickém panelu.

  \item \textbf{Pokrytí \ac{KV} vs.~Gini} ($r = -0{,}445$; $\rho = -0{,}286$):
    tento pár vykazuje nejmarkantnější divergenci mezi Pearsonem a~Spearmanem (rozdíl 0,159), která diagnosticky ukazuje na přítomnost outlierů zvyšujících Pearsonův koeficient.
    Nejpravděpodobnějším kandidátem je Francie: s~pokrytím \ac{KS} přes \SI{95}{\percent} (erga omnes) vykazuje relativně vysoké Gini (30--32), protože smluvní komprese mezd je přebita vysokou nerovností kapitálových příjmů a~vrcholných odměn.
    Průřezová hodnota $r_{\text{akt.}} = -0{,}238$ (\emph{slabá záporná}) dále ilustruje, proč nelze z~jednoduchého průřezu soudit o~teorii: v~roce 2024 komprimují Gini napříč EU27 silné transferové systémy a~progresivní zdanění, přičemž \ac{KV} je jen jedním z~více kolineárních institucio\-nálních faktorů.
    Panelový koeficient ($-0{,}445$) sledující \emph{změny} v~čase v~rámci zemí je proto informačně hodnotnější: dokumentuje, že kdy\-koli pokrytí \ac{KS} v~dané zemi kleslo, Gini se tendovalo zvýšit.

  \item \textbf{Pokrytí \ac{KV} vs.~příjem/\ac{HDP}} ($r = 0{,}227$; $\rho = 0{,}230$):
    \emph{slabá kladná} asociace, numericky konzistentní v~obou metodách.
    Relativně nízká hodnota je z~velké části artefaktem konstrukce ukazatele: \ac{HDP}/ob.~je sám o~sobě pozitivně korelován s~\ac{KV} pokrytím ($r = 0{,}471$), takže ve jmenovateli podílu čistý~příjem/\ac{HDP} oba členy rostou společně s~pokrytím a~efekt se vzájemně tlumí.
    Tento výsledek tedy nesvědčí o~tom, že \ac{KV} instituce nezvedají mzdy, ale o~tom, že tento konkrétní relativní ukazatel není citlivým testem dané hypotézy --- absolutní příjmová míra (předchozí pár) zůstává spolehlivějším indikátorem.

  \item \textbf{Hustota odborů vs.~Gini} ($r = -0{,}454$; $\rho = -0{,}455$):
    tento pár je statisticky nejrobustnějším v~celé tabulce: Pearsonovo $r$ a~Spearmanovo $\rho$ jsou prakticky totožné (rozdíl $< 0{,}001$), což svědčí o~přibližně lineárním vztahu bez systematického vlivu odlehlých hodnot.
    \emph{Střední záporná} asociace je stabilní v~celém panelu i~v~aktuálním průřezu ($r_{\text{akt.}} = -0{,}541$).
    Shoda obou koeficientů rovněž zpětně potvrzuje, že divergence u~páru \ac{KV}~vs.~Gini je specifická pro pokrytí (kde Francie tvoří strukturální outlier), nikoliv obecná vlastnost vztahu institucí \ac{KV} k~nerovnosti.
\end{itemize}

Žádná z~výše uvedených korelací nefalzifikuje teorii mzdové komprese prostřednictvím kolektivního vyjednávání.
Klíčovým argumentem je, že jednoduchou bivariátní korelací na heterogenním průřezu 27 zemí lze jen obtížně izolovat marginální příspěvek \ac{KV} pokrytí --- vzdor tomu jsou všechny korelace ve správném předpokládaném směru bez jediné výjimky.
Kauzální identifikaci přináší panelová analýza s~fixními efekty a~přirozené experimenty: metaanalytická literatura (např.~Jaumotte \& Osorio Buitron 2015; Card et al.~2004) konzistentně dokládá statisticky i~věcně významný efekt \ac{KV} institucí na mzdovou kompresi.
Korelace v~tabulce tedy plní funkci deskriptivního podkladu motivujícího institucionální případ\-ovou studii ČR --- nikoliv jako izolovaný kauzální důkaz, ale jako konzistentní mezinárodní vzorec, v~němž ČR se svým pokrytím $\approx 32\,\%$ systematicky vykazuje horšíu výsledky než srovnatelné ekonomiky s~vyšším pokrytím \ac{KS}.

Tabulka \ref{tab:coverage_correlation_table} shrnuje výsledky korelační analýzy provedené na párovém panelu zemí OECD/\ac{EU} za dostupná roky od roku 2004.
Zdrojem dat pro pokrytí kolektivního vyjednávání je databáze OECD/AIAS ICTWSS (ukazatel CBC, míra ERB), pro hustotu odborů tatáž databáze (TUD); výstupní proměnné pocházejí z~Eurostatu (lfsa\_ewhun2, nama\_10\_pc, ilc\_di12, earn\_nt\_net).
Každý řádek reportuje: počet párových pozorování $n$ (geo$\times$rok po párování datových sad), Pearsonovo $r$ na celém panelu (míra lineární asociace), Spearmanovo $\rho$ jako robustní neparametrická alternativa rezistentní vůči odlehlým hodnotám, a~$r_{\text{akt.}}$ = Pearsonovo $r$ pro průřezová data z~roku 2024.
Síla korelací je hodnocena dle Evansovy (1996) škály: $|r| < 0{,}20$ velmi slabá; $0{,}20$--$0{,}39$ slabá; $0{,}40$--$0{,}59$ střední; $0{,}60$--$0{,}79$ silná; $\geq 0{,}80$ velmi silná.

Výsledky jsou interpretovány po jednotlivých párech:

\begin{itemize}
  \item \textbf{Pokrytí \ac{KV} vs.~pracovní doba} ($r = -0{,}492$; $\rho = -0{,}384$):
    panelová Pearsonova korelace svědčí o~\emph{střední záporné} asociaci --- vyšší pokrytí \ac{KS} je systematicky spojeno s~kratší průměrnou pracovní dobou.
    Mírný pokles na Spearmanovo $\rho$ (slabá--střední) naznačuje přítomnost několika málo vlivných pozorování posouvajících Pearsonův koeficient, nikoliv výrazný outlier narušující charakter vztahu.
    Průřezová korelace $r_{\text{akt.}} = -0{,}536$ (\emph{střední záporná}) se oproti panelové hodnotě spíše zesílila --- vztah tedy není artefaktem historické variace, ale přetrvává i~v~aktuálním mezinárodním průřezu.

  \item \textbf{Pokrytí \ac{KV} vs.~\ac{HDP}/ob.~(\ac{PPS})} ($r = 0{,}471$; $\rho = 0{,}456$):
    \emph{střední kladná} asociace, vysoce konzistentní mezi Pearsonem a~Spearmanem --- rozdíl pouhých 0,015 indikuje přibližně lineární vztah bez dominantních odlehlých pozorování.
    Průřezová hodnota $r_{\text{akt.}} = 0{,}598$ se přibližuje hranici \emph{silné korelace} (Evans: 0,60), naznačujíce, že v~aktuálním průřezu roku 2024 je vazba mezi institucemi \ac{KV} a~životní úrovní ještě zřetelnější než v~historickém panelu.

  \item \textbf{Pokrytí \ac{KV} vs.~Gini} ($r = -0{,}445$; $\rho = -0{,}286$):
    tento pár vykazuje nejmarkantnější divergenci mezi Pearsonem a~Spearmanem (rozdíl 0,159), která diagnosticky ukazuje na přítomnost outlierů zvyšujících Pearsonův koeficient.
    Nejpravděpodobnějším kandidátem je Francie: s~pokrytím \ac{KS} přes \SI{95}{\percent} (erga omnes) vykazuje relativně vysoké Gini (30--32), protože smluvní komprese mezd je přebita vysokou nerovností kapitálových příjmů a~vrcholných odměn.
    Průřezová hodnota $r_{\text{akt.}} = -0{,}238$ (\emph{slabá záporná}) dále ilustruje, proč nelze z~jednoduchého průřezu soudit o~teorii: v~roce 2024 komprimují Gini napříč EU27 silné transferové systémy a~progresivní zdanění, přičemž \ac{KV} je jen jedním z~více kolineárních institucio\-nálních faktorů.
    Panelový koeficient ($-0{,}445$) sledující \emph{změny} v~čase v~rámci zemí je proto informačně hodnotnější: dokumentuje, že kdy\-koli pokrytí \ac{KS} v~dané zemi kleslo, Gini se tendovalo zvýšit.

  \item \textbf{Pokrytí \ac{KV} vs.~příjem/\ac{HDP}} ($r = 0{,}227$; $\rho = 0{,}230$):
    \emph{slabá kladná} asociace, numericky konzistentní v~obou metodách.
    Relativně nízká hodnota je z~velké části artefaktem konstrukce ukazatele: \ac{HDP}/ob.~je sám o~sobě pozitivně korelován s~\ac{KV} pokrytím ($r = 0{,}471$), takže ve jmenovateli podílu čistý~příjem/\ac{HDP} oba členy rostou společně s~pokrytím a~efekt se vzájemně tlumí.
    Tento výsledek tedy nesvědčí o~tom, že \ac{KV} instituce nezvedají mzdy, ale o~tom, že tento konkrétní relativní ukazatel není citlivým testem dané hypotézy --- absolutní příjmová míra (předchozí pár) zůstává spolehlivějším indikátorem.

  \item \textbf{Hustota odborů vs.~Gini} ($r = -0{,}454$; $\rho = -0{,}455$):
    tento pár je statisticky nejrobustnějším v~celé tabulce: Pearsonovo $r$ a~Spearmanovo $\rho$ jsou prakticky totožné (rozdíl $< 0{,}001$), což svědčí o~přibližně lineárním vztahu bez systematického vlivu odlehlých hodnot.
    \emph{Střední záporná} asociace je stabilní v~celém panelu i~v~aktuálním průřezu ($r_{\text{akt.}} = -0{,}541$).
    Shoda obou koeficientů rovněž zpětně potvrzuje, že divergence u~páru \ac{KV}~vs.~Gini je specifická pro pokrytí (kde Francie tvoří strukturální outlier), nikoliv obecná vlastnost vztahu institucí \ac{KV} k~nerovnosti.
\end{itemize}

Žádná z~výše uvedených korelací nefalzifikuje teorii mzdové komprese prostřednictvím kolektivního vyjednávání.
Klíčovým argumentem je, že jednoduchou bivariátní korelací na heterogenním průřezu 27 zemí lze jen obtížně izolovat marginální příspěvek \ac{KV} pokrytí --- vzdor tomu jsou všechny korelace ve správném předpokládaném směru bez jediné výjimky.
Kauzální identifikaci přináší panelová analýza s~fixními efekty a~přirozené experimenty: metaanalytická literatura (např.~Jaumotte \& Osorio Buitron 2015; Card et al.~2004) konzistentně dokládá statisticky i~věcně významný efekt \ac{KV} institucí na mzdovou kompresi.
Korelace v~tabulce tedy plní funkci deskriptivního podkladu motivujícího institucionální případ\-ovou studii ČR --- nikoliv jako izolovaný kauzální důkaz, ale jako konzistentní mezinárodní vzorec, v~němž ČR se svým pokrytím $\approx 32\,\%$ systematicky vykazuje horšíu výsledky než srovnatelné ekonomiky s~vyšším pokrytím \ac{KS}.

\input{texparts/python/coverage_correlation_table}



